% Options for packages loaded elsewhere
% Options for packages loaded elsewhere
\PassOptionsToPackage{unicode}{hyperref}
\PassOptionsToPackage{hyphens}{url}
\PassOptionsToPackage{dvipsnames,svgnames,x11names}{xcolor}
%
\documentclass[
  letterpaper,
  DIV=11,
  numbers=noendperiod]{scrreprt}
\usepackage{xcolor}
\usepackage{amsmath,amssymb}
\setcounter{secnumdepth}{5}
\usepackage{iftex}
\ifPDFTeX
  \usepackage[T1]{fontenc}
  \usepackage[utf8]{inputenc}
  \usepackage{textcomp} % provide euro and other symbols
\else % if luatex or xetex
  \usepackage{unicode-math} % this also loads fontspec
  \defaultfontfeatures{Scale=MatchLowercase}
  \defaultfontfeatures[\rmfamily]{Ligatures=TeX,Scale=1}
\fi
\usepackage{lmodern}
\ifPDFTeX\else
  % xetex/luatex font selection
\fi
% Use upquote if available, for straight quotes in verbatim environments
\IfFileExists{upquote.sty}{\usepackage{upquote}}{}
\IfFileExists{microtype.sty}{% use microtype if available
  \usepackage[]{microtype}
  \UseMicrotypeSet[protrusion]{basicmath} % disable protrusion for tt fonts
}{}
\makeatletter
\@ifundefined{KOMAClassName}{% if non-KOMA class
  \IfFileExists{parskip.sty}{%
    \usepackage{parskip}
  }{% else
    \setlength{\parindent}{0pt}
    \setlength{\parskip}{6pt plus 2pt minus 1pt}}
}{% if KOMA class
  \KOMAoptions{parskip=half}}
\makeatother
% Make \paragraph and \subparagraph free-standing
\makeatletter
\ifx\paragraph\undefined\else
  \let\oldparagraph\paragraph
  \renewcommand{\paragraph}{
    \@ifstar
      \xxxParagraphStar
      \xxxParagraphNoStar
  }
  \newcommand{\xxxParagraphStar}[1]{\oldparagraph*{#1}\mbox{}}
  \newcommand{\xxxParagraphNoStar}[1]{\oldparagraph{#1}\mbox{}}
\fi
\ifx\subparagraph\undefined\else
  \let\oldsubparagraph\subparagraph
  \renewcommand{\subparagraph}{
    \@ifstar
      \xxxSubParagraphStar
      \xxxSubParagraphNoStar
  }
  \newcommand{\xxxSubParagraphStar}[1]{\oldsubparagraph*{#1}\mbox{}}
  \newcommand{\xxxSubParagraphNoStar}[1]{\oldsubparagraph{#1}\mbox{}}
\fi
\makeatother

\usepackage{color}
\usepackage{fancyvrb}
\newcommand{\VerbBar}{|}
\newcommand{\VERB}{\Verb[commandchars=\\\{\}]}
\DefineVerbatimEnvironment{Highlighting}{Verbatim}{commandchars=\\\{\}}
% Add ',fontsize=\small' for more characters per line
\usepackage{framed}
\definecolor{shadecolor}{RGB}{241,243,245}
\newenvironment{Shaded}{\begin{snugshade}}{\end{snugshade}}
\newcommand{\AlertTok}[1]{\textcolor[rgb]{0.68,0.00,0.00}{#1}}
\newcommand{\AnnotationTok}[1]{\textcolor[rgb]{0.37,0.37,0.37}{#1}}
\newcommand{\AttributeTok}[1]{\textcolor[rgb]{0.40,0.45,0.13}{#1}}
\newcommand{\BaseNTok}[1]{\textcolor[rgb]{0.68,0.00,0.00}{#1}}
\newcommand{\BuiltInTok}[1]{\textcolor[rgb]{0.00,0.23,0.31}{#1}}
\newcommand{\CharTok}[1]{\textcolor[rgb]{0.13,0.47,0.30}{#1}}
\newcommand{\CommentTok}[1]{\textcolor[rgb]{0.37,0.37,0.37}{#1}}
\newcommand{\CommentVarTok}[1]{\textcolor[rgb]{0.37,0.37,0.37}{\textit{#1}}}
\newcommand{\ConstantTok}[1]{\textcolor[rgb]{0.56,0.35,0.01}{#1}}
\newcommand{\ControlFlowTok}[1]{\textcolor[rgb]{0.00,0.23,0.31}{\textbf{#1}}}
\newcommand{\DataTypeTok}[1]{\textcolor[rgb]{0.68,0.00,0.00}{#1}}
\newcommand{\DecValTok}[1]{\textcolor[rgb]{0.68,0.00,0.00}{#1}}
\newcommand{\DocumentationTok}[1]{\textcolor[rgb]{0.37,0.37,0.37}{\textit{#1}}}
\newcommand{\ErrorTok}[1]{\textcolor[rgb]{0.68,0.00,0.00}{#1}}
\newcommand{\ExtensionTok}[1]{\textcolor[rgb]{0.00,0.23,0.31}{#1}}
\newcommand{\FloatTok}[1]{\textcolor[rgb]{0.68,0.00,0.00}{#1}}
\newcommand{\FunctionTok}[1]{\textcolor[rgb]{0.28,0.35,0.67}{#1}}
\newcommand{\ImportTok}[1]{\textcolor[rgb]{0.00,0.46,0.62}{#1}}
\newcommand{\InformationTok}[1]{\textcolor[rgb]{0.37,0.37,0.37}{#1}}
\newcommand{\KeywordTok}[1]{\textcolor[rgb]{0.00,0.23,0.31}{\textbf{#1}}}
\newcommand{\NormalTok}[1]{\textcolor[rgb]{0.00,0.23,0.31}{#1}}
\newcommand{\OperatorTok}[1]{\textcolor[rgb]{0.37,0.37,0.37}{#1}}
\newcommand{\OtherTok}[1]{\textcolor[rgb]{0.00,0.23,0.31}{#1}}
\newcommand{\PreprocessorTok}[1]{\textcolor[rgb]{0.68,0.00,0.00}{#1}}
\newcommand{\RegionMarkerTok}[1]{\textcolor[rgb]{0.00,0.23,0.31}{#1}}
\newcommand{\SpecialCharTok}[1]{\textcolor[rgb]{0.37,0.37,0.37}{#1}}
\newcommand{\SpecialStringTok}[1]{\textcolor[rgb]{0.13,0.47,0.30}{#1}}
\newcommand{\StringTok}[1]{\textcolor[rgb]{0.13,0.47,0.30}{#1}}
\newcommand{\VariableTok}[1]{\textcolor[rgb]{0.07,0.07,0.07}{#1}}
\newcommand{\VerbatimStringTok}[1]{\textcolor[rgb]{0.13,0.47,0.30}{#1}}
\newcommand{\WarningTok}[1]{\textcolor[rgb]{0.37,0.37,0.37}{\textit{#1}}}

\usepackage{longtable,booktabs,array}
\newcounter{none} % for unnumbered tables
\usepackage{calc} % for calculating minipage widths
% Correct order of tables after \paragraph or \subparagraph
\usepackage{etoolbox}
\makeatletter
\patchcmd\longtable{\par}{\if@noskipsec\mbox{}\fi\par}{}{}
\makeatother
% Allow footnotes in longtable head/foot
\IfFileExists{footnotehyper.sty}{\usepackage{footnotehyper}}{\usepackage{footnote}}
\makesavenoteenv{longtable}
\usepackage{graphicx}
\makeatletter
\newsavebox\pandoc@box
\newcommand*\pandocbounded[1]{% scales image to fit in text height/width
  \sbox\pandoc@box{#1}%
  \Gscale@div\@tempa{\textheight}{\dimexpr\ht\pandoc@box+\dp\pandoc@box\relax}%
  \Gscale@div\@tempb{\linewidth}{\wd\pandoc@box}%
  \ifdim\@tempb\p@<\@tempa\p@\let\@tempa\@tempb\fi% select the smaller of both
  \ifdim\@tempa\p@<\p@\scalebox{\@tempa}{\usebox\pandoc@box}%
  \else\usebox{\pandoc@box}%
  \fi%
}
% Set default figure placement to htbp
\def\fps@figure{htbp}
\makeatother





\setlength{\emergencystretch}{3em} % prevent overfull lines

\providecommand{\tightlist}{%
  \setlength{\itemsep}{0pt}\setlength{\parskip}{0pt}}



 


\KOMAoption{captions}{tableheading}
\makeatletter
\@ifpackageloaded{bookmark}{}{\usepackage{bookmark}}
\makeatother
\makeatletter
\@ifpackageloaded{caption}{}{\usepackage{caption}}
\AtBeginDocument{%
\ifdefined\contentsname
  \renewcommand*\contentsname{Table of contents}
\else
  \newcommand\contentsname{Table of contents}
\fi
\ifdefined\listfigurename
  \renewcommand*\listfigurename{List of Figures}
\else
  \newcommand\listfigurename{List of Figures}
\fi
\ifdefined\listtablename
  \renewcommand*\listtablename{List of Tables}
\else
  \newcommand\listtablename{List of Tables}
\fi
\ifdefined\figurename
  \renewcommand*\figurename{Figure}
\else
  \newcommand\figurename{Figure}
\fi
\ifdefined\tablename
  \renewcommand*\tablename{Table}
\else
  \newcommand\tablename{Table}
\fi
}
\@ifpackageloaded{float}{}{\usepackage{float}}
\floatstyle{ruled}
\@ifundefined{c@chapter}{\newfloat{codelisting}{h}{lop}}{\newfloat{codelisting}{h}{lop}[chapter]}
\floatname{codelisting}{Listing}
\newcommand*\listoflistings{\listof{codelisting}{List of Listings}}
\makeatother
\makeatletter
\makeatother
\makeatletter
\@ifpackageloaded{caption}{}{\usepackage{caption}}
\@ifpackageloaded{subcaption}{}{\usepackage{subcaption}}
\makeatother
\usepackage{bookmark}
\IfFileExists{xurl.sty}{\usepackage{xurl}}{} % add URL line breaks if available
\urlstyle{same}
\hypersetup{
  pdftitle={kenai-chinook},
  pdfauthor={Benjamin Meyer},
  colorlinks=true,
  linkcolor={blue},
  filecolor={Maroon},
  citecolor={Blue},
  urlcolor={Blue},
  pdfcreator={LaTeX via pandoc}}


\title{kenai-chinook}
\author{Benjamin Meyer}
\date{}
\begin{document}
\maketitle

\renewcommand*\contentsname{Table of contents}
{
\setcounter{tocdepth}{2}
\tableofcontents
}

\bookmarksetup{startatroot}

\chapter*{Preface}\label{preface}
\addcontentsline{toc}{chapter}{Preface}

\markboth{Preface}{Preface}

This is a Quarto book.

To learn more about Quarto books visit
\url{https://quarto.org/docs/books}.

\begin{Shaded}
\begin{Highlighting}[]
\DecValTok{1} \SpecialCharTok{+} \DecValTok{1}
\end{Highlighting}
\end{Shaded}

\begin{verbatim}
[1] 2
\end{verbatim}

\bookmarksetup{startatroot}

\chapter{Introduction}\label{introduction}

This is a book created from markdown and executable code.

See @knuth84 for additional discussion of literate programming.

\begin{Shaded}
\begin{Highlighting}[]
\DecValTok{1} \SpecialCharTok{+} \DecValTok{1}
\end{Highlighting}
\end{Shaded}

\begin{verbatim}
[1] 2
\end{verbatim}

\bookmarksetup{startatroot}

\chapter{Summary}\label{summary}

In summary, this book has no content whatsoever.

\begin{Shaded}
\begin{Highlighting}[]
\DecValTok{1} \SpecialCharTok{+} \DecValTok{1}
\end{Highlighting}
\end{Shaded}

\begin{verbatim}
[1] 2
\end{verbatim}

\bookmarksetup{startatroot}

\chapter{Concept Note: An Integrated Doctoral Research Program on Kenai
River Chinook
Salmon}\label{concept-note-an-integrated-doctoral-research-program-on-kenai-river-chinook-salmon}

Connecting Habitat Condition, Water Quality, and Population Productivity

\hfill\break

\section{Overview}\label{overview}

Kenai River Chinook salmon (\emph{Oncorhynchus tshawytscha}) have
experienced prolonged abundance declines that have repeatedly triggered
emergency management actions, fishery closures, and, ultimately, a
federal commercial fishery disaster declaration for Upper Cook Inlet.
Despite decades of monitoring, the relative contributions of freshwater
habitat degradation, water quality impairment, marine survival shifts,
bycatch, and competition from hatchery-origin salmon have not been
rigorously quantified for this population. No integrated, multi-factor
research program exists to address these drivers in concert.

This document describes a proposed doctoral research program designed to
fill that gap. The program centers on the Kenai River and is organized
into four complementary research chapters, each targeted as a
standalone, peer-reviewed publication. The research will be conducted in
collaboration with Kenai Watershed Forum (KWF) and will draw on
long-term monitoring data from the Alaska Department of Fish and Game
(ADF\&G), NOAA, and the U.S. Geological Survey, supplemented by new
field data collection.

We are seeking funding support through the
\href{https://www.psmfc.org/procurement/upper-cook-inlet-salmon-disaster-research/}{PSMFC
Upper Cook Inlet Salmon Disaster Research program}, and we are actively
recruiting academic collaborators to serve as dissertation advisors and
committee members.

\begin{center}\rule{0.5\linewidth}{0.5pt}\end{center}

\section{Background}\label{background}

\subsection{Kenai River Chinook Salmon and the Upper Cook Inlet
Disaster}\label{kenai-river-chinook-salmon-and-the-upper-cook-inlet-disaster}

The Kenai River supports two distinct runs of Chinook salmon: an early
run returning primarily in May and June, and a late run returning in
July through September. The late run is the larger of the two and has
historically supported subsistence, sport, and commercial fisheries of
significant cultural and economic value throughout the Kenai Peninsula
and the broader Upper Cook Inlet.

Through the 2000s and 2010s, late-run Chinook abundance declined
substantially relative to historical levels. Spawner escapements
regularly fell below biological escapement goals, triggering Chinook
harvest restrictions on all user groups, as well full or partial
closures of commercial sockeye fisheries where Chinook was a non-target
species. By the early 2020s, cumulative losses had become severe enough
to warrant a federal commercial fishery disaster declaration, which
authorized dedicated research funding through NOAA and PSMFC.

The causes of this decline are not well characterized. A range of
factors have been implicated or suggested, including shifts in marine
survival associated with climate variability, competitive interactions
at sea with the expanding abundance of hatchery-origin pink and chum
salmon, freshwater habitat degradation driven by riparian development
and land use change along the lower Kenai River corridor, water quality
impairment related to sediment, temperature, and contaminants, and the
cumulative effects of harvest across multiple fisheries, including
bycatch in North Pacific trawl fisheries. No prior study has addressed
all of these factors simultaneously, or quantified their relative
contributions using a rigorous framework.

\subsection{The Research Gap}\label{the-research-gap}

Existing research on Kenai Chinook has generally addressed individual
drivers in isolation. Spawner-recruit analyses using standard regression
methods have been conducted periodically by ADF\&G but do not account
for observation error in escapement counts, do not integrate cohort
structure across age classes, and do not formally test hypotheses about
specific drivers. Habitat and water quality data have been collected by
KWF for many years but have not been formally linked to population
productivity metrics. Restoration activities, including riparian
planting and large wood placements, have been implemented but not
evaluated against measurable habitat or productivity outcomes.

The proposed doctoral program is designed to address these gaps through
coordinated, chapter-scale research that can collectively support
management decisions and advance the broader science of Pacific salmon
population dynamics.

\begin{center}\rule{0.5\linewidth}{0.5pt}\end{center}

\section{Proposed Research Program}\label{proposed-research-program}

The proposed doctoral program is organized around four research
chapters, described below. Each chapter is designed to stand alone as a
publishable unit, and each builds toward a synthesized understanding of
the factors shaping Kenai River Chinook salmon productivity. The program
is expected to span at least three to four years.

\subsection{Chapter 1: Long-Term Water Quality Monitoring and Trend
Analysis}\label{chapter-1-long-term-water-quality-monitoring-and-trend-analysis}

\textbf{Research question:} What are the long-term trends in water
quality parameters relevant to Chinook salmon in the Kenai River
watershed, and do those trends correlate with population productivity or
habitat condition indicators?

\textbf{Existing foundation:} KWF has maintained one of the most
comprehensive long-term water quality monitoring programs in Alaska,
with continuous records spanning 2000 to 2025 from sites distributed
across the Kenai River mainstem and key tributaries. A draft synthesis
report covering this 25-year record is currently in preparation (Meyer,
KWF; available online at
\url{https://kenai-watershed-forum.github.io/kenai-river-wqx/}), and
updates two prior published assessments of the same dataset (McCard
2007; Guerron Orejuela 2016). The draft report addresses 18 parameters,
including trace metals (arsenic, cadmium, chromium, copper, lead, zinc,
calcium, iron, magnesium), nutrients (nitrate, phosphorus), organic
compounds (BTEX), fecal coliform bacteria, and physical parameters (pH,
specific conductance, total suspended solids, turbidity, water
temperature).

\textbf{Approach:} This chapter will complete and formalize the
2000-2025 water quality synthesis as a peer-reviewed publication, with
statistical trend analysis characterizing temporal and spatial patterns
across the parameter suite. A second analytical component will link
water quality time series to independent indicators of Chinook salmon
productivity, testing whether any parameters exhibit temporal
correlations with escapement anomalies, juvenile condition metrics, or
the productivity index derived from Chapter 4.

\textbf{Significance:} No peer-reviewed synthesis of long-term water
quality trends in the Kenai River watershed currently exists. This
chapter will fill that gap, establish a citable baseline for future
management and litigation contexts, identify parameters of greatest
concern for salmon habitat, and connect KWF's monitoring investment
directly to salmon population outcomes.

\subsection{Chapter 2: Predictive Fish Habitat Mapping and Anadromous
Waters Catalog
Expansion}\label{chapter-2-predictive-fish-habitat-mapping-and-anadromous-waters-catalog-expansion}

\textbf{Research question:} Can predictive habitat modeling accurately
identify undocumented anadromous fish habitat across the Kenai Peninsula
Borough at landscape scale, and what does the resulting habitat map
reveal about the productive capacity available to Chinook salmon?

\textbf{Existing foundation:} In spring 2025, KWF was awarded a federal
grant through the National Coastal Resilience Fund to
\href{https://www.kenaiwatershed.org/news-media/fish-habitat-mapping-2025/}{explore
predictive habitat mapping for the Kenai Peninsula Borough}, a region
where fewer than 40\% of likely anadromous waters are currently
inventoried in the Alaska Department of Fish and Game's Anadromous
Waters Catalog (AWC). This initiative is modeled on work by the Hoonah
Native Forest Partnership in southeast Alaska, where a combination of
fieldwork, aerial imagery, and machine learning produced fish
presence/absence predictions with greater than 98\% accuracy within an
average of 200 feet. KWF is partnering with Romey Riverscape Sciences,
Terrainworks, and Cook Inletkeeper on the effort, with field surveys
underway beginning in 2025.

\textbf{Approach:} This chapter will develop the predictive habitat
model for the Kenai River watershed, validate it against field survey
data, and produce a publicly available habitat map that substantially
expands documented anadromous waters in the AWC. The analysis will
characterize predicted habitat distribution with respect to species
composition and life stage, with particular attention to Chinook salmon
rearing habitat in the lower Kenai River system. A synthesis section
will translate the modeled habitat estimates into productive capacity
metrics relevant to the spawner-recruit analysis in Chapter 4.

\textbf{Significance:} The existing AWC substantially underestimates
anadromous habitat across Alaska, leaving large areas of fish habitat
legally unprotected and effectively invisible to resource managers. For
the Kenai Peninsula Borough, completing a high-accuracy habitat
inventory would directly improve habitat conservation and development
permitting decisions, support restoration prioritization, and provide
the spatial context needed to interpret restoration outcomes (Chapter 3)
and population-level productivity trends (Chapter 4).

\subsection{Chapter 3: Corridor-Scale Habitat Quality Assessment and
Population-Level Productivity
Linkages}\label{chapter-3-corridor-scale-habitat-quality-assessment-and-population-level-productivity-linkages}

\textbf{Research question:} Across the full lower Kenai River corridor,
how much habitat meeting physical quality criteria for juvenile KRCS
rearing currently exists, how has that changed over time through
restoration, erosion, and development, and are cumulative changes in
physical habitat quality associated with detectable shifts in population
productivity?

\textbf{Relationship to existing work:} The Grant No.~26-084G project
(Tryon, Donnelly, and Muehlbauer, UAF/ADF\&G) is characterizing physical
habitat conditions and documenting juvenile Kenai River Chinook salmon
presence across seven bank cover types in the mainstem Kenai River. This
work will generate a valuable site-scale dataset on habitat conditions
and fish captures. Chapter 3 builds on those physical habitat
measurements while departing from an assumption embedded in the existing
project's framing: that fish presence or trap catch rate indicates
habitat quality.

Fish abundance in baited minnow traps is influenced by many factors
independent of habitat suitability, including differential trap
catchability across habitat types with different flow velocities and
structure, fish behavioral aggregation near certain features regardless
of habitat value, density-dependent displacement of subordinate fish
into suboptimal areas, and seasonal variation in activity. Using catch
rate as a proxy for quality risks confounding detectability with actual
habitat value. A more defensible approach evaluates habitat quality
directly through physical metrics that are mechanistically linked to
juvenile growth, condition, and survival.

\textbf{Approach:} Chapter 3 will use the physical habitat covariates
collected by the Grant No.~26-084G project (water depth, flow velocity,
substrate caliber, vegetation density, temperature, and bank cover type)
as the primary quality indicators, rather than fish presence as a proxy.
These covariates will be synthesized into a physical habitat quality
index grounded in existing literature on juvenile salmonid habitat
requirements. The index will then be applied at landscape scale using
remote sensing and historical aerial and satellite imagery to classify
bank condition across the full lower Kenai River corridor over a
multi-decade time series. The resulting corridor-scale time series of
physical habitat quality will be incorporated into the spawner-recruit
modeling framework from Chapter 4 to test whether cumulative changes in
habitat composition are associated with detectable shifts in population
productivity. This analysis will be coordinated with Tryon, Donnelly,
and Muehlbauer to ensure methodological consistency and avoid
duplication of field effort.

\textbf{Significance:} The Grant No.~26-084G project will produce
site-scale data on habitat conditions and fish catches. Chapter 3 asks
the next-order question: does the net change in physical habitat quality
across the whole corridor show up in the population record? Separating
habitat quality from fish detectability is essential to answering this
rigorously, and connecting physical habitat metrics to population
productivity provides the management-relevant result that restoration
funders and fisheries managers most need.

\subsection{Chapter 4: Spawner-Recruit Productivity Analysis with
Environmental and Anthropogenic
Covariates}\label{chapter-4-spawner-recruit-productivity-analysis-with-environmental-and-anthropogenic-covariates}

\textbf{Research question:} What are the relative contributions of
climate variability, North Pacific competition, bycatch, and density
dependence to the long-term productivity trend of Kenai River late-run
Chinook salmon?

\textbf{Approach:} This chapter aims to fit a Bayesian state-space
Ricker spawner-recruit model to reconstructed brood-year time series
derived from ADF\&G escapement, harvest, and age composition data,
following similar examples focused on Yukon River Chinook (DeFilippo et
al.~2026; Cunningham et al.~2018). The model will explicitly account for
the fact that spawner counts are imprecise measurements (sonar and weir
counts carry known error), that the underlying productivity of the
population may shift gradually over time for reasons not fully captured
by any single measured variable, and that fish removed by harvest and
bycatch in each year reduce the number of adults that ultimately return
to spawn. Confirmed available covariates include Gulf of Alaska sea
surface temperature and marine heatwave indices (NOAA ERSSTv5) and total
North Pacific pink and chum salmon abundance (Ruggerone and Irvine 2018
database). Two additional high-priority covariates are subject to
ongoing data availability assessments: (1) the separation of
hatchery-origin from wild pink and chum abundance, which would require
hatchery release records from ADF\&G and the North Pacific Anadromous
Fish Commission or genetic stock composition estimates; and (2) Gulf of
Alaska trawl bycatch apportioned to Kenai-origin Chinook, for which
observer-documented total bycatch exists but the depth and quality of
the Kenai-specific genetic apportionment record has not yet been
confirmed. A data inventory will be completed to determine which
covariates can be responsibly included, and any limitations will be
reported transparently. The model will use statistical techniques
designed to avoid overfitting - that is, to avoid drawing strong
conclusions from variables that may appear important by chance given the
limited length of available time series.

\textbf{Significance:} This chapter extends established methods from
Yukon Chinook modeling to a Cook Inlet context where a marine juvenile
abundance index is unavailable, and where bycatch attribution to a
specific population is particularly socially and politically salient. A
well-supported quantitative estimate of the relative magnitude of
bycatch effects, alongside climate and competition, would be a
significant contribution to ongoing management discussions.

\begin{center}\rule{0.5\linewidth}{0.5pt}\end{center}

\section{Partnerships and Institutional
Framework}\label{partnerships-and-institutional-framework}

The proposed program will be anchored at Kenai Watershed Forum, which
provides institutional knowledge, field infrastructure, long-term
datasets, and community relationships essential to each chapter.
Critically, KWF already has active, funded programs directly
corresponding to Chapters 1 and 2: the Kenai River baseline water
quality monitoring program (operational since 2000, with a 2000-2025
synthesis report in preparation) and the Kenai Peninsula Borough habitat
mapping initiative (funded by the National Coastal Resilience Fund,
begun 2025). The doctoral program would formalize and publish these
ongoing efforts as peer-reviewed research, while adding the
corridor-scale habitat quality assessment (Chapter 3) and productivity
modeling (Chapter 4) components that transform the collection of studies
into an integrated doctoral dissertation.

The doctoral student (or KWF employee, depending on the institutional
arrangement) will be supervised by a university-based advisory
committee.

We are actively seeking faculty members with expertise in one or more of
the following areas to serve as dissertation advisors or committee
members:

\begin{itemize}
\tightlist
\item
  Pacific salmon population dynamics and spawner-recruit modeling
\item
  Freshwater ecology and habitat assessment in Alaska
\item
  Bayesian statistical methods and state-space models
\item
  Watershed hydrology and geomorphology
\item
  Habitat restoration ecology and monitoring design
\end{itemize}

The logical home for the doctoral program is the University of Alaska
Fairbanks, though we welcome conversations with faculty at any
institution who share interest in the research questions described
above.

\begin{center}\rule{0.5\linewidth}{0.5pt}\end{center}

\section{Alignment with PSMFC Upper Cook Inlet Salmon Disaster Research
Priorities}\label{alignment-with-psmfc-upper-cook-inlet-salmon-disaster-research-priorities}

The PSMFC Upper Cook Inlet Salmon Disaster Research program specifically
targets research on Kenai late-run Chinook salmon and management of
Upper Cook Inlet District mixed-stock fisheries. The proposed doctoral
program directly addresses both priorities. Chapter 4 provides
quantitative information on the biological factors limiting
productivity, which is directly relevant to harvest management. Chapters
1, 2, and 3 address freshwater habitat conditions that affect productive
capacity and could inform long-term recovery strategies beyond harvest
management alone.

The program is designed to produce deliverables appropriate to the PSMFC
grant period: peer-reviewed publications, publicly available datasets
and code, and management-relevant technical summaries for ADF\&G and
other decision makers.

\begin{center}\rule{0.5\linewidth}{0.5pt}\end{center}

\section{Expected Outcomes}\label{expected-outcomes}

Upon completion, the doctoral program will deliver:

\begin{itemize}
\tightlist
\item
  Four peer-reviewed publications covering water quality trends, habitat
  extent modeling, restoration effectiveness, and productivity drivers.
\item
  Publicly available datasets and reproducible analysis code for all
  quantitative components.
\item
  Management-relevant summaries and briefings for ADF\&G, PSMFC, NOAA,
  and Kenai Peninsula stakeholders.
\item
  A spatially explicit habitat model and accessibility map for the Kenai
  River system.
\item
  A quantitative estimate of the relative contributions of climate,
  competition, bycatch, and density dependence to the Kenai late-run
  Chinook productivity time series.
\end{itemize}

\begin{center}\rule{0.5\linewidth}{0.5pt}\end{center}

\section{Budget and Timeline}\label{budget-and-timeline}

The program is expected to span at least three years, with a proposed
budget aligned with the PSMFC Upper Cook Inlet Salmon Disaster Research
program ceiling of approximately \$882,000. A detailed budget will be
developed in coordination with KWF and the supervising university, and
will reflect costs associated with personnel (graduate student or
postdoctoral researcher and KWF staff time), field data collection,
analytical computing, travel, publication costs, and indirect costs.

{\def\LTcaptype{none} % do not increment counter
\begin{longtable}[]{@{}
  >{\raggedright\arraybackslash}p{(\linewidth - 2\tabcolsep) * \real{0.5000}}
  >{\raggedright\arraybackslash}p{(\linewidth - 2\tabcolsep) * \real{0.5000}}@{}}
\toprule\noalign{}
\begin{minipage}[b]{\linewidth}\raggedright
Year
\end{minipage} & \begin{minipage}[b]{\linewidth}\raggedright
Primary Activities
\end{minipage} \\
\midrule\noalign{}
\endhead
\bottomrule\noalign{}
\endlastfoot
1 & Data compilation and quality control; Chapter 1 trend analysis;
Chapter 2 habitat model development; coordination with
Tryon/Donnelly/Muehlbauer on Chapter 3 physical habitat covariates \\
2 & Chapter 2 AWC expansion; Chapter 3 physical habitat index
development and remote sensing corridor classification; Chapter 4 cohort
reconstruction and model development; draft manuscripts for Chapters 1
and 2 \\
3 & Chapter 4 model fitting and covariate analysis; draft manuscripts
for Chapters 3 and 4; synthesis and dissemination \\
\end{longtable}
}

\begin{center}\rule{0.5\linewidth}{0.5pt}\end{center}

\section{Contact and Next Steps}\label{contact-and-next-steps}

We welcome inquiries from potential funders, academic advisors, and
collaborating researchers. For information about this proposal or to
discuss potential involvement, please contact:

\textbf{Benjamin Meyer}\\
Kenai Watershed Forum\\
ben@kenaiwatershed.org\\
\href{https://www.benjamin-meyer.net/}{www.benjamin-meyer.net}

We are preparing a formal pre-proposal for submission to the PSMFC Upper
Cook Inlet Salmon Disaster Research program upon release of the 2026
Request for Proposals. Interested faculty are encouraged to reach out as
early as possible, as institutional arrangements for the doctoral
program will need to be finalized prior to proposal submission.

\begin{center}\rule{0.5\linewidth}{0.5pt}\end{center}

\section{References}\label{references}

Cunningham, C. J., Westley, P. A. H., \& Adkison, M. D. (2018). Signals
of large-scale climate drivers, hatchery enhancement, and marine
survival in Yukon River Chinook salmon survival. \emph{Global Change
Biology}, 24(9), 4204-4217.

DeFilippo, L. B., et al.~(2026). Shifting stage-specific constraints on
productivity shape recovery potential of Chinook salmon.
\emph{Ecological Applications}.

\section{AI Disclosure}\label{ai-disclosure}

Some elements of this document were refined with the assistance of Posit
Assistant (Claude Sonnet 4.6), accessed via RStudio in April 2026. All
scientific content, framing decisions, and factual claims were reviewed
by the author.

\bookmarksetup{startatroot}

\chapter*{References}\label{references-1}
\addcontentsline{toc}{chapter}{References}

\markboth{References}{References}

\protect\phantomsection\label{refs}




\end{document}
